\section{Eksperimen dan Analisis}
\subsection{Tujuan dan Rancangan Eksperimen}
Eksperimen dirancang untuk membandingkan kualitas solusi, durasi, konsistensi,
serta perilaku konvergensi ketiga algoritma. Evaluasi GA memisahkan pengaruh
banyak generasi dan jumlah populasi melalui dua seri eksperimen terkontrol.

\subsubsection{Komposisi Dataset}
Dataset dipilih untuk mengevaluasi persoalan pemilihan dan penempatan paket.
Komposisi sebaiknya mencakup keterbatasan ruang, keterbatasan berat, serta
paket fragile. Rasio volume dan berat berikut dapat dilaporkan untuk menjelaskan
karakteristik kasus:
\begin{equation}
 R_V=\frac{\sum_i w_i^0\ell_i^0h_i^0}{WLH},\qquad
 R_M=\frac{\sum_i m_i}{C}.
\end{equation}
Rasio ini mendeskripsikan input, bukan bukti bahwa suatu susunan dapat dimuat.


\subsubsection{Pengujian Kebenaran Sebelum Eksperimen}
Sebelum hasil performa digunakan, pengujian harus memeriksa paket yang tepat
menyentuh dinding, paket yang memotong dinding, kapasitas tepat batas dan
berlebih, tabrakan volume, paket di lantai, paket melayang, support parsial,
fragile sebagai penopang, serta perpindahan masuk/keluar. Perubahan paket
penopang harus diikuti pemeriksaan paket yang terdampak.

Pengujian integrasi memeriksa bahwa final state yang ditampilkan memiliki
nilai objektif dan status validitas yang sama dengan output tersimpan.
Untuk GA, jumlah individu, statistik populasi, dan hasil decoding perlu
diperiksa; untuk SA, temperatur dan probabilitas penerimaan perlu konsisten.


\subsubsection{Prosedur Run dan Pencatatan}
Setiap varian hill-climbing yang dievaluasi dan SA dijalankan tiga kali pada
konfigurasi yang dinyatakan. Untuk GA disiapkan dua seri: tiga variasi iterasi
dengan populasi tetap, serta tiga variasi populasi dengan iterasi tetap.
Setiap konfigurasi pada setiap seri dijalankan tiga kali, sehingga rencana
ini menghasilkan 18 run GA per dataset.

Seed, parameter, lingkungan, dan definisi interval waktu dicatat. Durasi
pencarian sebaiknya dipisahkan dari rendering agar biaya visualisasi tidak
mengaburkan perbandingan. Initial state yang sebanding dapat dipakai untuk
hill-climbing dan SA; GA memakai populasi awal sehingga perbedaannya harus
dinyatakan. Perbandingan berdasarkan iterasi saja tidak menyamakan biaya
komputasi karena satu generasi GA mengevaluasi banyak individu.


\subsection{Hasil Eksperimen Hill-Climbing}


Untuk steepest-ascent, stochastic, dan sideways, jumlah iterasi hingga berhenti
dicatat. Sideways juga memerlukan parameter maksimum sideways move. Random
restart memerlukan maksimum restart, jumlah restart aktual, dan iterasi
per restart; rincian tersebut ditambahkan apabila varian ini dipilih.
% Visualisasi aktual dapat disisipkan setelah tersedia.

% Visualisasi aktual dapat disisipkan setelah tersedia.


\subsection{Hasil Eksperimen Simulated Annealing}


Frekuensi stuck hanya dapat ditafsirkan setelah definisi operasional dan
aturan penghitungannya ditetapkan. Grafik eksponensial perlu mencantumkan
konvensi $\Delta E$, penanganan temperatur nol, dan apakah data hanya
mencakup kandidat yang menurunkan nilai.
% Visualisasi aktual dapat disisipkan setelah tersedia.

% Visualisasi aktual dapat disisipkan setelah tersedia.

% Visualisasi aktual dapat disisipkan setelah tersedia.


\subsection{Hasil Eksperimen Genetic Algorithm}
\subsubsection{Pengaruh Banyak Iterasi}
Jumlah populasi ditetapkan sebesar $P_0$, sedangkan banyak generasi divariasikan
menjadi $G_1,G_2,G_3$ yang berbeda. Parameter lainnya dijaga tetap. Setiap
konfigurasi dijalankan tiga kali. Generasi digunakan sebagai istilah iterasi GA.

\subsubsection{Pengaruh Jumlah Populasi}
Banyak generasi ditetapkan sebesar $G_0$, sedangkan jumlah populasi
divariasikan menjadi $P_1,P_2,P_3$ yang berbeda. Setiap konfigurasi juga
dijalankan tiga kali.
\begin{table}[H]\centering\small
\caption{Rancangan matriks konfigurasi eksperimen GA}\label{tab:ga-config}
\begin{tabularx}{\linewidth}{@{}lYYYc@{}}\toprule
Seri & Konfigurasi & Populasi & Generasi & Run\\\midrule
Iterasi & I1 & $P_0$ & $G_1$ & 3\\
Iterasi & I2 & $P_0$ & $G_2$ & 3\\
Iterasi & I3 & $P_0$ & $G_3$ & 3\\
Populasi & P1 & $P_1$ & $G_0$ & 3\\
Populasi & P2 & $P_2$ & $G_0$ & 3\\
Populasi & P3 & $P_3$ & $G_0$ & 3\\\bottomrule
\end{tabularx}\end{table}


\subsubsection{Statistik dan Visualisasi Populasi}
Jika $\mathcal{P}_g$ merupakan populasi pada generasi $g$ dengan ukuran $P$,
statistik yang dicatat adalah
\begin{equation}
 f_{\max}(g)=\max_{s\in\mathcal{P}_g} f(s),\qquad
 \bar f(g)=\frac{1}{P}\sum_{s\in\mathcal{P}_g}f(s).
\end{equation}
Apabila fitness internal berbeda dari objektif, grafik wajib tetap diberi
label yang sesuai dan metode agregasinya dijelaskan. Statistik objektif
individu tidak valid tidak boleh dicampur tanpa penjelasan.
% Visualisasi aktual dapat disisipkan setelah tersedia.

% Visualisasi aktual dapat disisipkan setelah tersedia.

% Visualisasi aktual dapat disisipkan setelah tersedia.


\subsection{Analisis Perbandingan Algoritma}
\subsubsection{Kualitas Solusi dan Acuan Optimalitas}
Kualitas dibandingkan melalui nilai objektif pada state valid. Apabila
optimum $f^*$ diketahui dari pembuktian atau solver pembanding yang sah,
gap dapat dihitung sebagai $(f^*-f)/f^*$ untuk $f^*>0$.
Jika optimum tidak diketahui, gunakan istilah solusi terbaik yang ditemukan.
Untuk kasus dengan $v_i\geq0$, $U=\sum_i v_i$ merupakan batas atas sederhana,
tetapi dapat sangat longgar karena mengabaikan kendala ruang dan berat.
Selisih terhadap $U$ tidak boleh dinyatakan sebagai gap terhadap optimum pasti.


\subsubsection{Durasi, Konsistensi, dan Konvergensi}
Tiga run dapat diringkas menggunakan rata-rata, minimum, maksimum, dan
simpangan baku sampel nilai objektif serta durasi. Jumlah pengulangan yang
kecil membatasi kekuatan generalisasi. Grafik menjelaskan apakah kenaikan
nilai berlangsung cepat, berhenti pada plateau, atau berfluktuasi.


\subsubsection{Pengaruh Parameter Genetic Algorithm}
Analisis variasi generasi dan populasi dilakukan terpisah sesuai variabel
kontrol. Bahas perubahan kualitas akhir, keragaman hasil, dan durasi.
Peningkatan populasi atau generasi tidak diasumsikan selalu memberi hasil
lebih baik; kesimpulan harus mengikuti data yang diperoleh.


\subsection{Implementasi dan Evaluasi Bonus}
Rancangan bonus memanfaatkan atribut ETA yang telah tersedia pada paket
untuk mempertimbangkan urgensi pengiriman. Evaluasi bonus dipisahkan dari
evaluasi objektif wajib agar perubahan kriteria penilaian dapat dianalisis.
Spesifikasi juga menyediakan pilihan bonus seluruh varian hill-climbing dan
multi-truk; pemanfaatan ETA tidak otomatis berarti kedua fitur tersebut
ikut diterapkan~\cite{spesifikasi}.
